\documentclass{article}

\setlength{\headsep}{0.75 in}
\setlength{\parindent}{0 in}
\setlength{\parskip}{0.1 in}

%=====================================================
% Add PACKAGES Here (You typically would not need to):
%=====================================================

\usepackage{xcolor}
\usepackage[margin=1in]{geometry} % Used for page margins, often replaces fullpage
\usepackage{amsmath,amsthm,amssymb} % amssymb added for common math symbols

\usepackage{enumitem} % For custom list environments

\usepackage{tcolorbox}
% \usepackage{fullpage} % Geometry handles this better
\usepackage[parfill]{parskip} % For paragraph spacing
\usepackage{hyperref}
\usepackage[nameinlink, capitalise]{cleveref} % For smart referencing

\usepackage{algorithm2e} % For algorithms
\usepackage{tikz} % For drawing diagrams
\usepackage{siunitx} % For units
\usepackage{graphicx}

\usetikzlibrary{arrows.meta, positioning}

\crefname{algocf}{Algorithm}{Algorithms}

\RestyleAlgo{ruled} % Style for algorithms

\definecolor{linkc}{rgb}{0.1, 0.5, 0.7}
\definecolor{citec}{rgb}{0.6, 0.3, 0.7}
\definecolor{urlc}{rgb}{0.5, 0.1, 0.2}
\hypersetup{
    colorlinks=true,
    linkcolor=linkc,
    citecolor=citec,
    urlcolor=urlc
}

\renewcommand{\d}{\,\mathrm{d}} % Differential d


%=====================================================
% Ignore This Part (But Do NOT Delete It:)
%=====================================================

\theoremstyle{definition}
\newtheorem{problem}{Problem}
\newtheorem*{fun}{Fun with Algorithms}
\newtheorem*{challenge}{Challenge Yourself}
\def\fline{\rule{0.75\linewidth}{0.5pt}}
\newcommand{\finishline}{\begin{center}\fline\end{center}}
\newtheorem*{solution*}{Solution}
\newenvironment{solution}{\begin{solution*}}{{\finishline} \end{solution*}}
\newcommand{\grade}[1]{\hfill{\textbf{($\mathbf{#1}$ points)}}}
\newcommand{\thisdate}{\today}
\newcommand{\thissemester}{\textbf{Rutgers: Spring 2026}}
\newcommand{\thiscourse}{CS 344: Algorithms} 
\newcommand{\thishomework}{Number} 
\newcommand{\thisname}{Name} 
\newcommand{\thisextension}{Yes/No} 

\headheight 40pt              
\headsep 10pt


\newcommand{\thisheading}{
   \noindent
   \begin{center}
   \framebox{
      \vbox{\vspace{2mm}
    \hbox to 6.28in { \textbf{\thiscourse \hfill \thissemester} }
       \vspace{4mm}
       \hbox to 6.28in { {\Large \hfill Homework \#\thishomework \hfill} }
       \vspace{2mm}
         \hbox to 6.28in { { \hfill \hfill} }
       \vspace{2mm}
       \hbox to 6.28in { \emph{Name(s): \thisname \hfill }}
      \vspace{2mm}}
      }
   \end{center}
   \bigskip
}

%=====================================================
% Some useful MACROS (you can define your own in the same exact way also)
%=====================================================


\newcommand{\ceil}[1]{{\left\lceil{#1}\right\rceil}}
\newcommand{\floor}[1]{{\left\lfloor{#1}\right\rfloor}}
\newcommand{\prob}[1]{\Pr\paren{#1}}
\newcommand{\expect}[1]{\Exp\bracket{#1}}
\newcommand{\var}[1]{\textnormal{Var}\bracket{#1}}
\newcommand{\set}[1]{\ensuremath{\left\{ #1 \right\}}}
\newcommand{\poly}{\mbox{\rm poly}}


%=====================================================
% Fill Out This Part With Your Own Information:
%=====================================================


\renewcommand{\thishomework}{4} %Homework number
\renewcommand{\thisname}{FIRST$_1$ LAST$_1$, FIRST$_2$ LAST$_2$, FIRST$_3$ LAST$_3$} % Enter your name here


\begin{document}

\thisheading

\subsection*{Homework Policy}
\begin{itemize}
\item You may consult all course materials (lecture notes, textbook, slides, etc.) while writing your solutions. No external resources are permitted.


\item \textbf{All submissions should be typeset in LaTeX, handwritten solutions are not accepted}. 

\item Remember to always \textbf{prove the correctness} of your algorithms and \textbf{analyze their running time} (or any other efficiency measure asked in the question).

\item Homework assignments are completed and submitted in fixed groups of 3 students. Only \textbf{one member of the group submits} the full solution on Canvas. The submission must include the names of all group members. The other two members do \textbf{not need to submit anything separately}.

\item You may discuss the problems with any of your classmates, even those outside your group. However, \textbf{each group must write and submit their solutions independently}.

\item Attempt the challenge problems only if you find them interesting. 
\end{itemize}



\subsection*{Acknowledgment (to be reviewed and edited as needed)}
\begin{tcolorbox}[colframe=black,colback=white,arc=5mm]
I, \textbf{Iris Martinez (bam322)} worked individually on this assignment. \textbf{I hereby affirm that all external assistance utilized in the preparation of this document has been properly acknowledged}.
\end{tcolorbox}

\finishline
\newpage


\newpage


\begin{problem}
	Suppose we have $n$ persons with their heights given in an array $P[1:n]$ and $n$ skis with heights given in an array $S[1:n]$. Design and analyze a greedy algorithm that in $O(n\log{n})$ time 
	assign a ski to each person so that the average difference between the height of a person and their assigned ski is minimized. The algorithm should output a permutation $\sigma$ of $\set{1,\ldots,n}$ with the meaning 
	that person $i$ is assigned the ski $\sigma(i)$ such that 
	\[
		\frac1n \cdot \sum_{i=1}^{n} |P[i] - S[\sigma(i)]|
	\]
	is minimized. 
	
	\medskip

\paragraph{Example:} If $P=[3,2,5,1]$ and $B=[5,7,2,9]$, we can output $\sigma = [2,1,4,3]$ with value
\[
	\frac{1}{4} \cdot  \sum_{i=1}^{4} |P[i]-B[\sigma(i)]| = \frac14 \cdot (|P[1]-B[2]| + |P[2] - B[1]| + |P[3] - B[4]| + |P[4]-B[3]|) = \frac{1}{4} \cdot ({4+3+4+1}) = 3, 
\]
which you can also check is the minimum value. 


\end{problem} 





\begin{solution}

%=================================================
    % 1. THE ALGORITHM %=====================================================
    \

    \textbf{Algorithm:}
    \begin{itemize}
        \item sort each array. This takes nlogn time. Now take the shortest person and match them to the shortest ski. Do this over until all skis are assigned.
    \end{itemize}




%=================================================
    % 2. PROOF OF CORRECTNESS %=====================================================
    \

    \textbf{Proof: } This must give the lowest average gap between person height and ski height because, given an order of skis where the shortest person was not matched with the shortest ski, then the gap between the shortest person and their assigned ski would be greater, therefore increasing the average gap by at least the amount improved by using that ski to match with another person.
    Formally:
    Assume an optimal pairing of people and skis A with an average difference of b, where you could turn A into the greedy assignment G with one swap. In pairing A, the average difference is $\frac{P_1 - S_2 + P_2 - S_1 + ... + P_n - S_n}{n}$. If we swapped this to the greedy solution, for the greedy solution to be as good or better we must prove this inequality: $|P_1 - S_2| + |P_2 - S_1| \ge |P_1 - S_1| + |P_2 - S_2|$ where the heights of $P_1 < P_2$ and the height of $S_1 < S_2$. The improvement gained by assigning $S_1$ to $P_2$ will make the match between $S_2$ and $P_1$ worst by the same amount or less than the amount of the improvement gained by the swap. Therefore using the greedy ordering is always at least as good or the same as a some other ordering, and the greedy approach is proven.




%=================================================
    % 3. RUN TIME ANALYSIS %=====================================================
    
    \textbf{Runtime: } Sorting an array takes nlogn time with mergesort. Running that twice gives 2nlogn. The pairing takes n time. Therefore in the limit the algorithim is $O(nlogn)$


\end{solution}




\newpage




\begin{problem}
    Given a set of jobs, each with a deadline and a processing time, schedule the jobs to minimize the maximum lateness, where lateness is defined as the time difference between a job’s completion and its deadline (if it’s completed after the deadline). Write an algorithm and prove its correctness.
\end{problem}






\begin{solution}

    %====================%
    %= 1. THE ALGORITHM =%
    %====================%
    \

    \textbf{Algorithm: }
    \

    To minimize the max lateness of jobs, we sort the jobs by deadline and always pick the job in order of earliest deadline. 



%=================================================
    % 2. PROOF OF CORRECTNESS %=====================================================
    \

    \textbf{Proof: }
        Let's assume there was a solution other than the greedy solution. This would mean there exists an occasion where we are doing a job with a later deadline before one with an earlier deadline. In this case, inverting that execution order to match the greedy solution will either not increase the max lateness(in the case that the max lateness occurs due to another job). We also know that it might give a better solution by decreasing the lateness of the event with the earlier deadline. The event with the later deadline, even if it becomes more late, necessarily cannot increase our maximum because the new lateness must necessarily be less than the lateness that the job with the earlier deadline was facing in the "optimal" solution. Therefore the greedy solution is optimal.





%=================================================
    % 3. RUN TIME ANALYSIS %=====================================================
    \


    \textbf{Runtime: }
    This only needs to run one sort then retrieve the items in order of the sort, so it runs in $O(nlogn)$ time


\end{solution}


\newpage



\begin{problem}
    Alice wants to throw a party and she is trying to decide who to invite. She has $n$ people to choose from, and she knows which pairs of these people know each other. She wants to invite as many people as possible, subject to two constraints:

\begin{itemize}
    \item For each guest, there should be at least five other guests that they already know.
\item For each guest, there should be at least five other guests that they don't already know.

\end{itemize}
Describe and analyze an algorithm that computes the largest possible number of guests Alice can invite, given a list of $n$ people and the list of pairs who know each other.
\end{problem}





\begin{solution}

%=================================================
    % 1. THE ALGORITHM %=====================================================
    \

    \textbf{Algorithm: }

    The setup of a list of people and a list of pairs that know each other can represent a graph where every node is a person and every edge is a pair of people that know each other. 

    Iterate through all nodes, calculating the number of outbound edges for every node and the number of "unknown" nodes as well (which is the same as the total number of nodes N - the number of outgoing edges).

    For any node we find where either the number of outbound edges or unknown edges is less than five, we add them to a queue for deletion,D.
    
    For every node in the deletion queue, find and remove it from the graph. Then recalculate outgoing edges for all nodes connected to the node. Recalculate unkown nodes for all not related nodes. If any nodes break the rules, add them to the deltion queue. Repeat until the deltion queue is empty. The number of nodes in the graph is the maximum number of invitable people.

%=================================================
    % 2. PROOF OF CORRECTNESS %=====================================================
    \

    \textbf{Proof: }
    This algorithm must be correct because it progressively trims people that don't meet the conditions until it arrives at a maximal set of them to invite. If there were a party larger than one given by this algorithm, this means the algorithm must remove someone wrongly. Because our algorithm runs the outgoing edges and unknown nodes calculation on each node, it cannot ever remove a node that isn't below the thresholds. Therefore, the algorithm cannot arrive at a set that is too small. We can prove this by assuming that the optimal party, O, has one more node than ours, P. Lets call the "mistakenly" removed node of P, x. At the end of the loop where x was added to the deletion queue, every other node met the conditions, and thus $P \cap x$




%=================================================
    % 3. RUN TIME ANALYSIS %=====================================================
    \

    \textbf{Runtime: }
    The algorithm will run in $O(N^2)$ time. It will take $O(N^2)$ time to calculate the outgoing edges and unknown nodes values. In the worst case, we would cascade into a case where every single node would need to be removed, resulting in $N$ calculations for all $N$ nodes, giving a runtime of $O(2N^2)$ which simplifies to $O(N^2)$


\end{solution}

\newpage




\begin{problem}
A \textbf{walk} in a directed graph $G=(V,E)$ from a vertex $s$ to a vertex $t$, is a sequence of vertices $v_1,v_2,\ldots,v_k$ where $v_1 = s$ and $v_k=t$ such that for any $i < k$,  $(v_i,v_{i+1})$ is an edge in $G$. The \emph{length} of a walk
is defined as the number of vertices inside it minus one, i.e., the number of edges (so the walk $v_1,v_2,\ldots,v_k$ has length $k-1$). 

Note that the only difference 
of a walk with a \textbf{path} we defined in the course is that a walk can contain the same vertex (or edge) more than once, while a path consists of only distinct vertices and edges. 

Design and analyze an $O(n+m)$ time algorithm that given a directed graph $G=(V,E)$ and two vertices $s$ and $t$, outputs \emph{Yes} if there is a walk from $s$ to $t$ in $G$ \underline{whose length is even}, and \emph{No} otherwise. 


\end{problem}


\begin{solution}

%=================================================
    % 1. THE ALGORITHM %=====================================================





%=================================================
    % 2. PROOF OF CORRECTNESS %=====================================================





%=================================================
    % 3. RUN TIME ANALYSIS %=====================================================


\end{solution}

\newpage





\begin{problem}
	We have an undirected graph $G=(V,E)$ and two vertices $s$ and $t$. Additionally, we are given a list of \emph{new} edges $F$ which do \emph{not} belong to the graph $G$. For any edge $f \in F$, 
	we define the graph $G_f$ as the graph obtained by adding the edge $f$ to the original graph $G$. 
	
	Design an algorithm that in time $O(n+m+|{F}|)$ finds the edge $f \in F$ such that 
	the distance from $s$ to $t$ in $G_f$ is minimized among all choices of an edge from $F$.  
\end{problem}



\begin{solution}

%=================================================
    % 1. THE ALGORITHM %=====================================================





%=================================================
    % 2. PROOF OF CORRECTNESS %=====================================================





%=================================================
    % 3. RUN TIME ANALYSIS %=====================================================


\end{solution}

\newpage




\begin{problem}
    Consider a directed graph $G$, where each edge is colored either red, white, or blue. A walk in $G$ is called a French flag walk if its sequence of edge colors is red, white, blue, red, white, blue, and so on. More formally, a walk $v_0 \rightarrow v_1 \rightarrow \cdots \rightarrow v_k$ is a French flag walk if, for every integer $i$, the edge $v_i \rightarrow v_{i+1}$ is red if $i \bmod 3=0$, white if $i \bmod 3=1$, and blue if $i \bmod 3=2$.

Describe an algorithm to find all vertices in $G$ that can be reached from a given vertex $v$ through a French flag walk. 
\end{problem}



\begin{solution}

%=================================================
    % 1. THE ALGORITHM %=====================================================





%=================================================
    % 2. PROOF OF CORRECTNESS %=====================================================





%=================================================
    % 3. RUN TIME ANALYSIS %=====================================================


\end{solution}

\newpage




\begin{problem}
Suppose we are given both an undirected graph $G$ with weighted edges and a minimum spanning tree $T$ of $G$.
\begin{enumerate}
    \item [(a)]Describe an algorithm to update the minimum spanning tree when the weight of a single edge $e$ is decreased.
\item[(b)] Describe an algorithm to update the minimum spanning tree when the weight of a single edge $e$ is increased.

\end{enumerate}

\end{problem}



\begin{solution}

%=================================================
    % 1. THE ALGORITHM %=====================================================





%=================================================
    % 2. PROOF OF CORRECTNESS %=====================================================





%=================================================
    % 3. RUN TIME ANALYSIS %=====================================================


\end{solution}

\newpage




\finishline

\begin{challenge}
Suppose you are given an arbitrary directed graph $G$ in which each edge is colored either red or blue, along with two special vertices $s$ and $t$.
\begin{enumerate}
    \item[(a)] Describe an algorithm that either computes a walk from $s$ to $t$ such that the pattern of red and blue edges along the walk is a palindrome, or correctly reports that no such walk exists.
\item[(b)] Describe an algorithm that either computes the shortest walk from $s$ to $t$ such that the pattern of red and blue edges along the walk is a palindrome, or correctly reports that no such walk exists.
	\end{enumerate}
\end{challenge}



\begin{solution}

%=================================================
    % 1. THE ALGORITHM %=====================================================





%=================================================
    % 2. PROOF OF CORRECTNESS %=====================================================





%=================================================
    % 3. RUN TIME ANALYSIS %=====================================================


\end{solution}

\newpage





\end{document}
